\documentclass[11pt]{article}
\usepackage[margin=1in]{geometry}
\usepackage{amsmath,amssymb,amsthm}
\usepackage[hidelinks]{hyperref}

\newtheorem{theorem}{Theorem}
\newtheorem{lemma}[theorem]{Lemma}

\title{Proof of the Line-Metric Distortion Conjecture with Preference Intensities}
\author{Salem Levinstein}
\date{November 2026}

\begin{document}
\maketitle

\begin{abstract}
We prove the line-metric distortion conjecture of M.~Abbaszadeh, A.~Ansarifar,
M.~Latifian, and M.~Seddighin (arXiv:2601.02095v1, Section~6). For
two-alternative elections on the $1$-Euclidean line under mandatory
preference-intensity elicitation with threshold $\alpha \in (0, 1]$, we prove
that the rule $f^{\mathsf{tal}}_\alpha$ has worst-case metric distortion
exactly
\[
M(\alpha) := \max\!\left(\frac{3 - \alpha}{1 + \alpha},\, 1 + 2\alpha\right),
\]
and that no deterministic voting rule achieves strictly smaller distortion.
The proof linearizes the worst-case distortion condition into two linear forms
per alternative and shows that every cross-alternative pair of forms has
nonpositive coefficients at $t = M(\alpha)$. The entire argument is formally
verified in Lean~4.
\end{abstract}

\section{Introduction}

In metric distortion with preference intensities~\cite{AALS2026}, voters and
alternatives lie in a metric space $(X, d)$, and each voter reports ordinal
preferences augmented with an intensity threshold $\alpha \in (0, 1]$. For two
alternatives $a_1, a_2$ under mandatory elicitation, each voter $i \in [N]$
reports one of four types:
\begin{align*}
a_1 \succ a_2 &: \alpha\, d(i, a_2) < d(i, a_1) \le d(i, a_2), &
a_1 \gg a_2 &: d(i, a_1) \le \alpha\, d(i, a_2), \\
a_2 \succ a_1 &: \alpha\, d(i, a_1) < d(i, a_2) \le d(i, a_1), &
a_2 \gg a_1 &: d(i, a_2) \le \alpha\, d(i, a_1).
\end{align*}
Let $E$ be an $N$-voter profile with type counts $n = (n_1, n_2, n_3, n_4)$
corresponding to $(a_1 \succ a_2,\, a_2 \succ a_1,\, a_1 \gg a_2,\,
a_2 \gg a_1)$. On the $1$-Euclidean line $\mathbb{R}$, a placement
$(a_1, a_2, x) \in \mathbb{R} \times \mathbb{R} \times \mathbb{R}^N$ is
\emph{consistent} with $E$ if every voter $i$ satisfies the inequality of its
reported type with $d(i, a) = |x_i - a|$. The social cost of $a \in \{a_1, a_2\}$
is $\mathrm{SC}_x(a) := \sum_{i=1}^N |x_i - a|$. For $c \in \{a_1, a_2\}$, the
worst-case distortion of $c$ on $E$ is
\[
\mathrm{dist}_\alpha(c, E) := \inf \bigl\{ t \ge 1 : \forall \text{ consistent }
(a_1, a_2, x),\; \mathrm{SC}_x(c) \le t \min(\mathrm{SC}_x(a_1), \mathrm{SC}_x(a_2)) \bigr\}
\in [1, \infty].
\]
The rule $f^{\mathsf{tal}}_\alpha$ selects the alternative minimizing
$\mathrm{dist}_\alpha(c, E)$, and the distortion of a deterministic rule $f$ is
$\mathrm{dist}_\alpha(f) := \sup_E \mathrm{dist}_\alpha(f(E), E)$ over all
nonempty profiles $E$.

In Section~6 of~\cite{AALS2026}, Abbaszadeh, Ansarifar, Latifian, and
Seddighin supported the following conjecture by numerical simulations on all
profiles with $N = 1000$:
\begin{quote}
Considering elections with two alternatives and assuming that the metric is
limited to the $1$-Euclidean space, we have
$\mathrm{dist}_\alpha(f^{\mathsf{tal}}_\alpha) = \max\!\left(\frac{3 - \alpha}{1 + \alpha},\, 2\alpha + 1\right)$.
\end{quote}

We settle this conjecture affirmatively and show that the bound is tight for
every deterministic voting rule.

\section{Main Theorem and Proof}

\begin{theorem}\label{thm:main}
Fix $\alpha \in (0, 1]$ and set
$M(\alpha) := \max\bigl(\frac{3 - \alpha}{1 + \alpha},\, 1 + 2\alpha\bigr)$.
For two-alternative elections on $\mathbb{R}$ under mandatory elicitation:
\begin{enumerate}
\item Every profile $E$ satisfies
$\min(\mathrm{dist}_\alpha(a_1, E), \mathrm{dist}_\alpha(a_2, E)) \le M(\alpha)$.
\item Every deterministic voting rule $f$ satisfies
$\mathrm{dist}_\alpha(f) \ge M(\alpha)$.
\end{enumerate}
In particular, $\mathrm{dist}_\alpha(f^{\mathsf{tal}}_\alpha) = M(\alpha)$.
\end{theorem}

We begin with the per-voter upper bounds after affine normalization to
$(a_1, a_2) = (0, 1)$. For $t \ge 1$, set
\[
g_t(y) := |y| - t|y - 1|, \qquad
u_t := \frac{1 - t\alpha}{1 + \alpha}, \qquad
v_t := \frac{1 - t\alpha}{1 - \alpha} \quad (\alpha < 1).
\]

\begin{lemma}\label{lem:per-agent}
Let $\alpha \in (0, 1]$, $t \ge 1$, and $y \in \mathbb{R}$.
\begin{enumerate}
\item If $y$ satisfies $a_1 \succ a_2$ at $(0, 1)$, then
$g_t(y) \le \frac{1 - t}{2}$.
\item If $y$ satisfies $a_1 \gg a_2$ at $(0, 1)$, then
$g_t(y) \le \frac{\alpha - t}{1 + \alpha}$.
\item If $y$ satisfies $a_2 \gg a_1$ at $(0, 1)$, then $g_t(y) \le 1$.
\item If $\alpha < 1$ and $y$ satisfies $a_2 \succ a_1$ at $(0, 1)$, then
$g_t(y) \le \max(u_t, v_t)$. If $\alpha = 1$, neither $a_1 \succ a_2$ nor
$a_2 \succ a_1$ is satisfiable.
\end{enumerate}
\end{lemma}

\begin{proof}
(1) From $|y| \le |y - 1|$ and the triangle inequality
$1 \le |y| + |y - 1| \le 2|y - 1|$, we obtain $|y - 1| \ge \frac{1}{2}$. Since
$1 - t \le 0$, we have $g_t(y) \le (1 - t)|y - 1| \le \frac{1 - t}{2}$.

(2) From $|y| \le \alpha|y - 1|$ and $1 \le |y| + |y - 1| \le (1 + \alpha)|y - 1|$,
we have $|y - 1| \ge \frac{1}{1 + \alpha}$. Since $\alpha - t \le 0$,
$g_t(y) \le (\alpha - t)|y - 1| \le \frac{\alpha - t}{1 + \alpha}$.

(3) Since $|y| \le |y - 1| + 1$ and $t \ge 1$,
$g_t(y) \le |y| - |y - 1| \le 1$.

(4) When $\alpha = 1$, the condition $1 \cdot |y| < |y - 1| \le |y|$ is
contradictory (and likewise for $a_1 \succ a_2$). Now assume $\alpha < 1$ and
$\alpha|y| < |y - 1| \le |y|$. If $y \le 1$, then $|y - 1| \le |y|$ forces
$y \ge 0$, so $\alpha y < 1 - y$, giving $y < \frac{1}{1 + \alpha}$. Since
$1 + t > 0$,
\[
g_t(y) = y - t(1 - y) = (1 + t)y - t \le \frac{1 + t}{1 + \alpha} - t = u_t.
\]
If $y > 1$, then $\alpha y < y - 1$, giving $y > \frac{1}{1 - \alpha}$. Since
$1 - t \le 0$,
\[
g_t(y) = y - t(y - 1) = (1 - t)y + t \le \frac{1 - t}{1 - \alpha} + t = v_t.
\]
In both cases $g_t(y) \le \max(u_t, v_t)$.
\end{proof}

\begin{proof}[Proof of Theorem~\ref{thm:main}]
First observe that for any $t \ge 1$, because social costs are nonnegative, the
inequality $\mathrm{SC}_x(a_1) \le t\min(\mathrm{SC}_x(a_1), \mathrm{SC}_x(a_2))$
holds if and only if $\mathrm{SC}_x(a_1) \le t\,\mathrm{SC}_x(a_2)$. If
$a_1 = a_2$, this is immediate. Otherwise $s := a_2 - a_1 \ne 0$, and rescaling
$y_i := (x_i - a_1)/s$ maps $(a_1, a_2)$ to $(0, 1)$, preserves every
preference type, and scales both social costs by $|s|^{-1}$. Hence
$\mathrm{dist}_\alpha(a_1, E) \le t$ holds whenever
$\sum_{i=1}^N g_t(y_i) \le 0$ for all consistent normalized placements $y$.

\medskip
\noindent\textbf{Part 1 (Upper bound).}
Write $M := M(\alpha) \ge 1 + 2\alpha > 1$.

\emph{Case $\alpha < 1$.} For $n = (n_1, n_2, n_3, n_4) \in \mathbb{R}_{\ge 0}^4$,
define the linear forms
\[
B_t^{(1)}(n) := \frac{1 - t}{2} n_1 + u_t n_2 + \frac{\alpha - t}{1 + \alpha} n_3 + n_4,
\qquad
B_t^{(2)}(n) := \frac{1 - t}{2} n_1 + v_t n_2 + \frac{\alpha - t}{1 + \alpha} n_3 + n_4,
\]
and let $\sigma(n) := (n_2, n_1, n_4, n_3)$. By Lemma~\ref{lem:per-agent},
$\sum_{i=1}^N g_t(y_i) \le \max\bigl(B_t^{(1)}(n), B_t^{(2)}(n)\bigr)$ for
every consistent normalized placement $y$. Hence
$B_t^{(1)}(n) \le 0 \land B_t^{(2)}(n) \le 0$ implies
$\mathrm{dist}_\alpha(a_1, E) \le t$, and by swapping alternative labels,
$B_t^{(1)}(\sigma n) \le 0 \land B_t^{(2)}(\sigma n) \le 0$ implies
$\mathrm{dist}_\alpha(a_2, E) \le t$.

At $t = M$, consider the coefficients of $B_M^{(i)}(n) + B_M^{(j)}(\sigma n)$
for any $i, j \in \{1, 2\}$:
\begin{itemize}
\item On $n_3$ and $n_4$:
$\frac{\alpha - M}{1 + \alpha} + 1 = \frac{1 + 2\alpha - M}{1 + \alpha} \le 0$,
because $M \ge 1 + 2\alpha$.
\item On $n_1$ and $n_2$ with $v_M$:
$\frac{1 - M}{2} + v_M = \frac{(3 - \alpha) - M(1 + \alpha)}{2(1 - \alpha)} \le 0$,
because $M \ge \frac{3 - \alpha}{1 + \alpha}$.
\item On $n_1$ and $n_2$ with $u_M$:
$\frac{1 - M}{2} + u_M = \frac{(3 + \alpha) - M(1 + 3\alpha)}{2(1 + \alpha)} \le 0$,
because
\[
(3 - \alpha)(1 + 3\alpha) - (3 + \alpha)(1 + \alpha) = 4\alpha(1 - \alpha) > 0,
\]
which implies $M \ge \frac{3 - \alpha}{1 + \alpha} > \frac{3 + \alpha}{1 + 3\alpha}$.
\end{itemize}
Since $n_1, n_2, n_3, n_4 \ge 0$, every cross-sum
$B_M^{(i)}(n) + B_M^{(j)}(\sigma n)$ is nonpositive. Therefore we cannot have
both $\max(B_M^{(1)}(n), B_M^{(2)}(n)) > 0$ and
$\max(B_M^{(1)}(\sigma n), B_M^{(2)}(\sigma n)) > 0$. Consequently, at least
one of $a_1, a_2$ has distortion at most $M$.

\emph{Case $\alpha = 1$.} Here $M(1) = 3$. If any voter has type
$a_1 \succ a_2$ or $a_2 \succ a_1$, no consistent placement exists by
Lemma~\ref{lem:per-agent}(4) and the bound holds vacuously. Otherwise
$n_1 = n_2 = 0$, and Lemma~\ref{lem:per-agent} yields
$\sum_i g_3(y_i) \le -n_3 + n_4$ for $a_1$ and $-n_4 + n_3$ for $a_2$; at least
one of these two quantities is nonpositive.

\medskip
\noindent\textbf{Part 2 (Lower bound and tightness).}
Consider two symmetric two-voter profiles:
\begin{enumerate}
\item $E_{\mathrm{int}} := (a_1 \gg a_2,\, a_2 \gg a_1)$, for any
$\alpha \in (0, 1]$. Taking $(a_1, a_2) = (0, 1)$ and
$x = \bigl(\frac{\alpha}{1 + \alpha}, 1\bigr)$ gives a consistent placement
with $\mathrm{SC}_x(0) = \frac{1 + 2\alpha}{1 + \alpha}$ and
$\mathrm{SC}_x(1) = \frac{1}{1 + \alpha}$. Thus
$\mathrm{dist}_\alpha(a_1, E_{\mathrm{int}}) \ge 1 + 2\alpha$; by symmetry,
$\mathrm{dist}_\alpha(a_2, E_{\mathrm{int}}) \ge 1 + 2\alpha$ (and in fact
equality holds by Lemma~\ref{lem:per-agent}).
\item $E_{\mathrm{norm}} := (a_1 \succ a_2,\, a_2 \succ a_1)$, for
$\alpha \in (0, 1)$. Given any $t < \frac{3 - \alpha}{1 + \alpha}$, set
\[
\delta := \frac{(3 - \alpha) - t(1 + \alpha)}{2(1 - \alpha)} > 0, \qquad
\varepsilon := \frac{\delta}{2(|t| + 1)} > 0.
\]
Place $(a_1, a_2) = (0, 1)$ and
$x = \bigl(\frac{1}{2}, \frac{1}{1 - \alpha} + \varepsilon\bigr)$. Voter~1
satisfies $\frac{\alpha}{2} < \frac{1}{2} = \frac{1}{2}$, and voter~2 satisfies
$\alpha\bigl(\frac{1}{1 - \alpha} + \varepsilon\bigr) < \frac{\alpha}{1 - \alpha} + \varepsilon < \frac{1}{1 - \alpha} + \varepsilon$,
so $x$ is consistent. Moreover
$\mathrm{SC}_x(0) = \frac{3 - \alpha}{2(1 - \alpha)} + \varepsilon >
\frac{1 + \alpha}{2(1 - \alpha)} + \varepsilon = \mathrm{SC}_x(1)$, and
\[
\mathrm{SC}_x(0) - t\,\mathrm{SC}_x(1) = \delta + (1 - t)\varepsilon
\ge \delta - (|t| + 1)\varepsilon = \frac{\delta}{2} > 0.
\]
Hence $\mathrm{dist}_\alpha(a_1, E_{\mathrm{norm}}) \ge \frac{3 - \alpha}{1 + \alpha}$,
and by symmetry $\mathrm{dist}_\alpha(a_2, E_{\mathrm{norm}}) \ge \frac{3 - \alpha}{1 + \alpha}$
(with equality by Lemma~\ref{lem:per-agent}).
\end{enumerate}
At $\alpha = 1$, we have $\frac{3 - 1}{1 + 1} = 1 < 3 = 1 + 2\cdot 1$, so
$E_{\mathrm{int}}$ already witnesses $M(1) = 3$. Therefore every deterministic
rule $f$ satisfies $\mathrm{dist}_\alpha(f) \ge M(\alpha)$, and combining with
Part~1 gives $\mathrm{dist}_\alpha(f^{\mathsf{tal}}_\alpha) = M(\alpha)$.
\end{proof}

\section{Formal Verification}

The complete proof of Theorem~\ref{thm:main}---including affine normalization,
per-agent bounds, the cross-alternative pairing dichotomy, and both tightness
constructions---has been formally verified in Lean~4. In particular, universal
rule optimality (\texttt{any\_rule\_ruleDist\_ge}) and the exact distortion
equality (\texttt{fTal\_ruleDist\_eq}) are machine-checked using only the
standard foundational axioms (\texttt{propext}, \texttt{Classical.choice},
\texttt{Quot.sound}), with zero errors and zero warnings.

\begin{thebibliography}{9}
\bibitem{AALS2026}
M.~Abbaszadeh, A.~Ansarifar, M.~Latifian, and M.~Seddighin,
\emph{Metric Distortion with Preference Intensities},
arXiv:2601.02095v1 (2026), AAAI~2026.
\end{thebibliography}

\end{document}
